%!TEX program = xelatex
\documentclass[a4paper,11pt]{article}

\usepackage{kotex}
\usepackage[margin=2.5cm, top=3cm, bottom=3cm]{geometry}
\usepackage{xcolor}
\usepackage[most]{tcolorbox}
\usepackage{enumitem}
\usepackage{booktabs}
\usepackage{titlesec}
\usepackage{fancyhdr}
\usepackage{hyperref}
\usepackage{fontspec}
\usepackage{multirow}

% ── 색상 정의 ──────────────────────────────────────────────────────
\definecolor{primary}{RGB}{37, 99, 235}
\definecolor{green}{RGB}{22, 163, 74}
\definecolor{amber}{RGB}{217, 119, 6}
\definecolor{red}{RGB}{220, 38, 38}
\definecolor{indigo}{RGB}{79, 70, 229}
\definecolor{lightblue}{RGB}{239, 246, 255}
\definecolor{lightgreen}{RGB}{240, 253, 244}
\definecolor{lightamber}{RGB}{255, 251, 235}
\definecolor{lightred}{RGB}{254, 242, 242}
\definecolor{lightgray}{RGB}{249, 250, 251}

% ── 섹션 스타일 ────────────────────────────────────────────────────
\titleformat{\section}
  {\Large\bfseries\color{primary}}
  {\thesection.}{0.5em}{}[\vspace{-0.3em}\color{primary}\rule{\linewidth}{1pt}]

\titleformat{\subsection}
  {\normalsize\bfseries}
  {\thesubsection}{0.5em}{}

\titleformat{\subsubsection}
  {\small\bfseries\color{indigo}}
  {\thesubsubsection}{0.5em}{}

% ── 박스 스타일 ────────────────────────────────────────────────────
\tcbset{
  tipbox/.style={
    colback=lightblue, colframe=primary, fonttitle=\bfseries\small,
    title={도움말}, arc=3pt, boxrule=0.5pt, left=6pt, right=6pt, top=5pt, bottom=5pt
  },
  warnbox/.style={
    colback=lightamber, colframe=amber, fonttitle=\bfseries\small,
    title={주의}, arc=3pt, boxrule=0.5pt, left=6pt, right=6pt, top=5pt, bottom=5pt
  },
  stepbox/.style={
    colback=lightgreen, colframe=green, arc=3pt, boxrule=0.5pt,
    left=6pt, right=6pt, top=5pt, bottom=5pt
  },
  infobox/.style={
    colback=lightgray, colframe=gray!50, arc=3pt, boxrule=0.5pt,
    left=6pt, right=6pt, top=5pt, bottom=5pt
  }
}

% ── 헤더/푸터 ──────────────────────────────────────────────────────
\pagestyle{fancy}
\fancyhf{}
\fancyhead[L]{\small\color{gray} 양현고 수학여행 앱}
\fancyhead[R]{\small\color{gray} 교사 사용자 매뉴얼}
\fancyfoot[C]{\small\thepage}
\renewcommand{\headrulewidth}{0.4pt}

\setlist[enumerate]{leftmargin=1.5em, itemsep=2pt}
\setlist[itemize]{leftmargin=1.5em, itemsep=2pt, label={\textcolor{primary}{$\bullet$}}}

% ──────────────────────────────────────────────────────────────────
\begin{document}

% ── 표지 ──────────────────────────────────────────────────────────
\begin{titlepage}
  \centering
  \vspace*{3cm}
  {\Huge\bfseries\color{primary} 양현고 수학여행 앱}\par\vspace{0.5cm}
  {\LARGE 교사 사용자 매뉴얼}\par\vspace{1cm}
  \textcolor{gray}{\rule{8cm}{0.5pt}}\par\vspace{1cm}
  {\large 2026 수학여행 현장 관리 시스템}\par\vspace{0.5cm}
  {\normalsize\color{gray} v1.0 · 2026년 5월}\par
  \vfill
  \begin{tcolorbox}[colback=lightblue, colframe=primary, arc=6pt,
    boxrule=1pt, width=10cm]
    \centering
    \textbf{접속 주소}\par\vspace{0.3em}
    {\large\color{primary}\texttt{yanghyeon-tour-2026.vercel.app}}
  \end{tcolorbox}
\end{titlepage}

\tableofcontents
\newpage

% ──────────────────────────────────────────────────────────────────
\section{시작하기}

\subsection{교사 계정 로그인}

\begin{enumerate}
  \item 브라우저에서 \textbf{\texttt{yanghyeon-tour-2026.vercel.app}} 접속
  \item 로그인 화면에서 교사 이메일 및 비밀번호 입력
  \item 로그인 후 자동으로 \textbf{교사 포털}로 이동
\end{enumerate}

\begin{tcolorbox}[tipbox]
교사 계정이 없는 경우 관리자에게 교사 가입 코드를 받아 회원가입하세요. 교사 포털 접근에는 별도 코드 인증이 필요합니다.
\end{tcolorbox}

\subsection{교사 포털 홈 구성}

\begin{center}
\begin{tabular}{@{}lll@{}}
  \toprule
  \textbf{영역} & \textbf{내용} & \textbf{조건} \\
  \midrule
  \textbf{세션 배너} & 현재 진행 중인 점호로 빠른 이동 & 점호 세션 진행 시 표시 \\
  \textbf{진행 중 점호} & 세션 목록·진행률 확인 & 점호 세션 진행 시 표시 \\
  \textbf{핵심 액션} & 자주 쓰는 메뉴 바로가기 & 항상 표시 \\
  \textbf{조회 메뉴} & 숙소 배정, 일정, 연락처 & 항상 표시 \\
  \bottomrule
\end{tabular}
\end{center}

\subsection{담임 기본값 자동 설정}

담임 정보(담임학년·담임반)가 등록된 교사는 각 화면 접속 시 자신의 반이 자동으로 필터됩니다.

\begin{tcolorbox}[infobox]
담임 정보가 설정되지 않았거나 틀릴 경우 관리자에게 \textbf{담임 정보 동기화}를 요청하세요.
\end{tcolorbox}

% ──────────────────────────────────────────────────────────────────
\section{세션 배너 및 점호 진행 상황}

\subsection{세션 배너}

진행 중인 점호 세션이 있을 때 홈 화면 최상단에 색상 배너가 표시됩니다.

\begin{center}
\begin{tabular}{@{}cll@{}}
  \toprule
  \textbf{색상} & \textbf{점호 유형} & \textbf{이동 화면} \\
  \midrule
  파랑 & 승차점호 진행 중 & 승차 현황 모니터링 \\
  초록 & 정시점호 진행 중 & 야간 점호 현황 \\
  \bottomrule
\end{tabular}
\end{center}

배너를 탭하면 해당 점호 현황 화면으로 즉시 이동합니다.

\subsection{점호 세션 카드}

홈 화면의 \textbf{「진행 중인 점호」} 섹션에서 각 세션의 상세 현황을 확인할 수 있습니다.

\begin{itemize}
  \item \textbf{도넛 차트}: 전체 대비 점호 완료 비율 시각화
  \item \textbf{진행 막대}: 완료율 색상 표시 (빨강→주황→파랑→초록)
  \item \textbf{「점호 현황」 버튼}: 해당 세션 상세 화면으로 이동
\end{itemize}

\begin{tcolorbox}[warnbox]
점호 연장(+10분) 및 종료 버튼은 \textbf{관리자만} 사용할 수 있습니다.
\end{tcolorbox}

% ──────────────────────────────────────────────────────────────────
\section{승차점호 모니터링}

홈 화면 \textbf{「승차 현황」} 카드 또는 배너를 탭하면 승차점호 모니터링 화면으로 이동합니다.

\subsection{반별 그리드 (전체 현황)}

모든 반의 승차 현황이 \textbf{2열 카드 그리드}로 표시됩니다.

\begin{center}
\begin{tabular}{@{}cll@{}}
  \toprule
  \textbf{카드 색상} & \textbf{의미} & \textbf{상태} \\
  \midrule
  초록 테두리 & 전원 승차 완료 & 출발 가능 \\
  주황 테두리 & 일부 승차 & 진행 중 \\
  빨강 테두리 & 아직 미확인 & 점호 미시작 \\
  \bottomrule
\end{tabular}
\end{center}

\begin{itemize}
  \item \textbf{담임반}은 파란색 테두리(링)로 강조 표시됩니다.
  \item 상단 배너: 전체 반 중 완료된 반 수 / 전체 인원 대비 확인 인원 표시
  \item 모든 반이 완료되면 \textbf{「전체 승차 완료! 출발 가능」} 메시지 표시
\end{itemize}

\subsection{반 선택 후 버스 좌석 배치도}

카드를 탭하면 해당 반의 \textbf{버스 좌석 배치도}가 표시됩니다.

\begin{tcolorbox}[infobox]
\textbf{버스 좌석 배치 구조:}
\begin{itemize}
  \item 상단: 운전석
  \item 좌측 2열 + 통로 + 우측 2열 (4열 배치)
  \item 각 좌석: 학생 이름, 학년·반·번호 표시
  \item \textcolor{green}{\textbf{초록 좌석}}: 탑승 완료 \quad \textbf{흰 좌석}: 미탑승
\end{itemize}
\end{tcolorbox}

\begin{enumerate}
  \item 좌석을 \textbf{탭}하면 수동으로 탑승/취소 처리 가능
  \item 미탑승 학생은 화면 아래 \textbf{연락처 섹션}에 표시됨
  \item 연락처 버튼을 눌러 \textbf{학생 또는 보호자에게 즉시 전화}
\end{enumerate}

\begin{tcolorbox}[warnbox]
점호 시간이 만료된 세션에서는 수동 탑승 처리가 불가합니다.
\end{tcolorbox}

\subsection{호차별 좌석 뷰}

상단의 \textbf{「호차별」} 토글을 선택하면 버스 호차 전체를 한 화면에서 볼 수 있습니다.

% ──────────────────────────────────────────────────────────────────
\section{야간 점호 (정시점호) 모니터링}

\textbf{「점호 현황」} 버튼 또는 세션 배너를 통해 이동합니다.

\subsection{화면 구성}

\subsubsection{상단 헤더}

\begin{center}
\begin{tabular}{@{}ll@{}}
  \toprule
  \textbf{요소} & \textbf{설명} \\
  \midrule
  도넛 차트 & 담임반 전체 점호 완료율 \\
  전체 / 완료 / 미완 카운터 & 실시간 인원 수 \\
  이름·번호·호실 검색창 & 특정 학생 빠른 찾기 \\
  학년 필터 (전체/1/2/3학년) & 학년별 보기 \\
  반 필터 (전체/1\,--\,10반) & 반별 보기 \\
  뷰 토글 (리스트 / 방) & 표시 형태 전환 \\
  \bottomrule
\end{tabular}
\end{center}

\subsubsection{리스트 뷰}

학생 1명씩 행으로 표시되며, 상태 필터(전체/완료/미완료)로 필터링 가능합니다.

\begin{itemize}
  \item \textbf{초록 행}: 점호 완료 학생
  \item \textbf{흰 행}: 미완료 학생
  \item \textbf{「직접 점호」 버튼}: 학생을 직접 탭하여 점호 처리
  \item \textbf{「완료」 버튼}: 탭하면 취소 확인 후 점호 취소 가능
\end{itemize}

\subsubsection{방 뷰 (객실 평면도)}

각 \textbf{호실을 위에서 내려다본 평면도}로 표시합니다.

\begin{tcolorbox}[infobox]
\textbf{객실 평면도 구조:}
\begin{itemize}
  \item 굵은 테두리: 방의 벽
  \item 좌우 침대열: 헤드보드가 벽 쪽을 향한 침대 배치
  \item 중앙 점선: 통로 구분
  \item 하단: 출입구 표시
  \item \textcolor{green}{\textbf{초록 침대}}: 점호 완료 \quad \textbf{흰 침대}: 미완료
\end{itemize}
\end{tcolorbox}

\begin{itemize}
  \item 침대를 \textbf{탭}하면 해당 학생 점호/취소 처리
  \item 호실 헤더의 도넛 차트로 호실별 완료율 한눈에 파악
  \item 완료 호실: 초록 테두리 / 미완료: 주황 또는 빨강 테두리
\end{itemize}

\begin{tcolorbox}[tipbox]
정시점호 세션에서는 \textbf{방 뷰가 기본}으로 열립니다. 상단 토글로 리스트 뷰로 전환할 수 있습니다.
\end{tcolorbox}

\subsection{직접 점호 처리 절차}

\begin{tcolorbox}[stepbox]
\begin{enumerate}
  \item 리스트 뷰: 해당 학생 행 오른쪽 \textbf{「직접 점호」} 버튼 탭
  \item 방 뷰: 해당 학생의 \textbf{침대 칸} 탭
  \item 점호 완료 확인 토스트 메시지 확인
  \item 실수로 점호한 경우: 완료 상태 버튼을 다시 탭 → 취소 확인 → 취소 완료
\end{enumerate}
\end{tcolorbox}

% ──────────────────────────────────────────────────────────────────
\section{학생 검색}

홈 화면 \textbf{「학생 검색」} 카드를 탭하여 이동합니다.

\subsection{검색 및 필터}

\begin{itemize}
  \item 이름, 학년, 반, 번호, 호실로 \textbf{실시간 검색} 가능
  \item \textbf{학년 / 반 필터}: 담임반이 자동으로 기본 선택됨
  \item \textbf{「다른 반·전체 학년 보기」} 링크: 필터 초기화하여 전체 조회
\end{itemize}

\subsection{학생 상세 정보}

검색 결과에서 학생을 탭하면 상세 정보가 표시됩니다.

\begin{center}
\begin{tabular}{@{}ll@{}}
  \toprule
  \textbf{항목} & \textbf{내용} \\
  \midrule
  학년·반·번호, 이름 & 기본 신상 \\
  호차, 호실, 비행편 & 여행 배정 정보 \\
  학생 연락처 & 탭하면 즉시 전화 연결 \\
  보호자 연락처 & 탭하면 즉시 전화 연결 \\
  건강 요주의사항 & 의료 주의 정보 \\
  특이사항 & 기타 특기 사항 \\
  가입 상태 & 앱 가입 여부 확인 \\
  \bottomrule
\end{tabular}
\end{center}

\begin{tcolorbox}[tipbox]
연락처 옆 전화 아이콘을 탭하면 스마트폰 전화 앱으로 즉시 연결됩니다.
\end{tcolorbox}

% ──────────────────────────────────────────────────────────────────
\section{숙소 배정 조회}

홈 화면 \textbf{「숙소 배정」} 카드를 탭하여 이동합니다.

\begin{itemize}
  \item 담임반 학생들의 \textbf{호실 배정 현황}을 확인할 수 있습니다.
  \item \textbf{호실 카드}를 탭하면 해당 호실 배정 학생 목록이 펼쳐집니다.
  \item 각 학생의 \textbf{학생 연락처} 및 \textbf{보호자 연락처}에서 즉시 전화 가능
\end{itemize}

\begin{tcolorbox}[infobox]
호실 정보는 Excel 업로드 시 입력된 학생 배정 데이터를 기반으로 표시됩니다.
담임반 정보가 없는 경우 전체 학생이 표시됩니다.
\end{tcolorbox}

% ──────────────────────────────────────────────────────────────────
\section{사건사고 등록}

홈 화면 \textbf{「사건사고 등록」} 카드를 탭하여 이동합니다.

\subsection{사건사고 기록 방법}

\begin{tcolorbox}[stepbox]
\begin{enumerate}
  \item 「+ 새 사건사고」 버튼 탭
  \item \textbf{유형} 선택 (부상, 분실물, 질병, 기타 등)
  \item \textbf{심각도} 선택 (경미 / 중요 / 긴급)
  \item \textbf{관련 학생} 검색 후 선택
  \item \textbf{내용} 및 \textbf{조치사항} 입력
  \item \textbf{「등록」} 버튼 탭
\end{enumerate}
\end{tcolorbox}

\begin{tcolorbox}[warnbox]
\textbf{긴급} 등급의 사건은 즉시 관리자에게도 알려주세요. 앱의 사건 등록은 기록 목적이며, 응급 상황 시에는 119 신고 및 관리자 직접 연락이 우선입니다.
\end{tcolorbox}

\subsection{사건 목록 확인}

\begin{itemize}
  \item 등록된 사건 목록에서 \textbf{처리 여부} 확인 가능
  \item \textbf{「종결」} 처리: 사건이 해결되면 종결 표시
  \item 미처리 사건은 빨간색으로 강조 표시
\end{itemize}

% ──────────────────────────────────────────────────────────────────
\section{채팅 (공지방)}

홈 화면 \textbf{「채팅」} 카드를 탭하여 이동합니다.

\begin{itemize}
  \item \textbf{교직원 채팅방}: 교사 간 및 관리자와 소통하는 채팅방
  \item \textbf{학급 채팅방}: 담임 반 학생들에게 공지를 보내는 일방향 채팅방
  \item 학생들은 읽기만 가능하며, 교사·관리자만 메시지를 작성할 수 있습니다.
\end{itemize}

% ──────────────────────────────────────────────────────────────────
\section{여행 일정 / 비상 연락처}

\subsection{여행 일정}

\begin{itemize}
  \item \textbf{1\,--\,4일차 탭}으로 일차별 일정 확인
  \item 각 일정: 시작 시각, 일정명, 장소, 종료 시각, 비고 표시
  \item \textbf{「간략히 / 자세히」} 토글로 상세 정보 표시 조절
  \item 점호가 연결된 일정: 색상 배지 (초록: 정시점호 / 파랑: 승차점호)
\end{itemize}

\subsection{비상 연락처}

\begin{itemize}
  \item 구분(병원, 경찰서 등)별 그룹화된 연락처 목록 확인
  \item 전화 버튼으로 즉시 연결 가능
\end{itemize}

% ──────────────────────────────────────────────────────────────────
\section{자주 묻는 질문}

\begin{description}[leftmargin=0pt, style=nextline]
  \item[\textbf{Q. 담임반이 자동으로 설정되지 않습니다.}]
    관리자에게 담임 정보 동기화를 요청하세요. Excel 업로드 후 「담임 정보 동기화」를 실행해야 적용됩니다.

  \item[\textbf{Q. 학생이 점호를 못 했습니다.}]
    세션이 만료 전이라면 「직접 점호」 버튼으로 교사가 수동 처리할 수 있습니다. 세션 만료 후에는 관리자에게 세션 연장을 요청하세요.

  \item[\textbf{Q. 승차점호에서 좌석이 표시되지 않습니다.}]
    진행 중인 승차점호 세션이 없거나, 해당 반 학생 데이터가 없는 것입니다. 관리자에게 세션 시작을 요청하세요.

  \item[\textbf{Q. 직접 점호한 것을 취소해야 합니다.}]
    점호 완료 상태(초록색)의 버튼 또는 침대를 다시 탭하면 취소 확인 후 취소됩니다.

  \item[\textbf{Q. 학생 비밀번호를 재설정해야 합니다.}]
    학생 검색 화면에서 해당 학생을 검색한 후 계정 섹션의 「비밀번호 재설정 이메일 발송」 버튼을 사용하거나, 관리자에게 UID 초기화를 요청하세요.
\end{description}

\end{document}
